% Options for packages loaded elsewhere
\PassOptionsToPackage{unicode}{hyperref}
\PassOptionsToPackage{hyphens}{url}
\documentclass[
  a4paper,
]{ltjsarticle}
\usepackage{xcolor}
\usepackage[margin=25mm]{geometry}
\usepackage{amsmath,amssymb}
\setcounter{secnumdepth}{-\maxdimen} % remove section numbering
\usepackage{iftex}
\ifPDFTeX
  \usepackage[T1]{fontenc}
  \usepackage[utf8]{inputenc}
  \usepackage{textcomp} % provide euro and other symbols
\else % if luatex or xetex
  \usepackage{unicode-math} % this also loads fontspec
  \defaultfontfeatures{Scale=MatchLowercase}
  \defaultfontfeatures[\rmfamily]{Ligatures=TeX,Scale=1}
\fi
\usepackage{lmodern}
\ifPDFTeX\else
  % xetex/luatex font selection
\fi
% Use upquote if available, for straight quotes in verbatim environments
\IfFileExists{upquote.sty}{\usepackage{upquote}}{}
\IfFileExists{microtype.sty}{% use microtype if available
  \usepackage[]{microtype}
  \UseMicrotypeSet[protrusion]{basicmath} % disable protrusion for tt fonts
}{}
\makeatletter
\@ifundefined{KOMAClassName}{% if non-KOMA class
  \IfFileExists{parskip.sty}{%
    \usepackage{parskip}
  }{% else
    \setlength{\parindent}{0pt}
    \setlength{\parskip}{6pt plus 2pt minus 1pt}}
}{% if KOMA class
  \KOMAoptions{parskip=half}}
\makeatother
\usepackage{longtable,booktabs,array}
\newcounter{none} % for unnumbered tables
\usepackage{calc} % for calculating minipage widths
% Correct order of tables after \paragraph or \subparagraph
\usepackage{etoolbox}
\makeatletter
\patchcmd\longtable{\par}{\if@noskipsec\mbox{}\fi\par}{}{}
\makeatother
% Allow footnotes in longtable head/foot
\IfFileExists{footnotehyper.sty}{\usepackage{footnotehyper}}{\usepackage{footnote}}
\makesavenoteenv{longtable}
\usepackage{graphicx}
\makeatletter
\newsavebox\pandoc@box
\newcommand*\pandocbounded[1]{% scales image to fit in text height/width
  \sbox\pandoc@box{#1}%
  \Gscale@div\@tempa{\textheight}{\dimexpr\ht\pandoc@box+\dp\pandoc@box\relax}%
  \Gscale@div\@tempb{\linewidth}{\wd\pandoc@box}%
  \ifdim\@tempb\p@<\@tempa\p@\let\@tempa\@tempb\fi% select the smaller of both
  \ifdim\@tempa\p@<\p@\scalebox{\@tempa}{\usebox\pandoc@box}%
  \else\usebox{\pandoc@box}%
  \fi%
}
% Set default figure placement to htbp
\def\fps@figure{htbp}
\makeatother
\setlength{\emergencystretch}{3em} % prevent overfull lines
\providecommand{\tightlist}{%
  \setlength{\itemsep}{0pt}\setlength{\parskip}{0pt}}
\usepackage{bookmark}
\IfFileExists{xurl.sty}{\usepackage{xurl}}{} % add URL line breaks if available
\urlstyle{same}
\hypersetup{
  pdftitle={第四思考実験：固定作用 S から状態依存作用 S\_n へ ― 第三思考実験で外部計算していた相対論的重力軌道を、X\_\{n+1\}=S\_nX\_n の逐次相互作用から再生成できるか ―},
  pdfauthor={Noriaki Kihara (ORCID 0009-0004-6753-4020)},
  hidelinks,
  pdfcreator={LaTeX via pandoc}}

\title{第四思考実験：固定作用 \(S\) から状態依存作用 \(S_n\) へ ―
第三思考実験で外部計算していた相対論的重力軌道を、\(X_{n+1}=S_nX_n\)
の逐次相互作用から再生成できるか ―}
\author{Noriaki Kihara (ORCID 0009-0004-6753-4020)}
\date{2026-09-23 / v1.0 / Version DOI 10.5281/zenodo.22919788}

\sloppy
\emergencystretch=3em
\begin{document}
\maketitle

\textbf{著者:} 木原範昭 (Noriaki Kihara)\\
\textbf{ORCID:} 0009-0004-6753-4020\\
\textbf{Version DOI:} 10.5281/zenodo.22919788\\
\textbf{Concept DOI:} 10.5281/zenodo.22919787\\
\textbf{版:} v1.0\\
\textbf{日付:} 2026-09-23

\begin{center}\rule{0.5\linewidth}{0.5pt}\end{center}

\subsection{要旨}\label{ux8981ux65e8}

本稿は、第二思考実験と第三思考実験の間に残った一つの不整合から出発する。第二思考実験では、意味を定めない二つの値からなる内部状態
\(X_n\) に固定作用 \(S\)
を繰り返し作用させ、積の読み出しと面積の時計から逆二乗型の軌道を生成した。一方、第三思考実験では方向を逆にし、既知のニュートン力学およびポストニュートン近似によって二体軌道を外部で計算したうえで、その観測データから質量比・相対距離・相対時間を読み出せるかを検証した。この方法では読み出し能力は確認できたが、最も重要な軌道生成そのものが既知の重力方程式に残っていた。

そこで本稿では、第三思考実験で外部から与えていた相対論的軌道を、再び内部状態の逐次相互作用から生成できるかを調べた。固定した同一作用
\(S\)
のままでは、今回の状態表現と読み出しを保ったまま、軌道位置に依存する近点移動と放射反作用による軌道要素の変化を同時に表すための処理依存性を作用の外側に持ち込む必要が生じる。そこで最小の一般化として、作用自身を現在の関係状態の関数とし、

\[
X_{n+1}=S_nX_n,
\qquad
S_n=\mathcal S(\mathcal R_n),
\qquad
\det S_n=1
\]

とした。ここで \(n\) は物理時刻ではなく処理番号であり、物理時間 \(t\)
は状態列から別途読み出す。

保存部分には Schwarzschild 束縛軌道の Darwin
パラメータ表示を用い、その角度進行を任意次数の漸化式で構成した。さらに第三思考実験と同じ
2.5PN 放射反作用を、軌道要素 \(p_n,e_n\)
の局所更新として加えた。最初の強場条件では、放射反作用を追加することで基準PN軌道に対する規格化半径
RMS 誤差が 2.51\% から 1.95\% に改善した。

次に、\(S_n\) の生成則を変更せず、\(p=24.3,60,120\) と
\(e=0.20,0.45,0.55\) の9条件へ適用した。9ケースの規格化半径 RMS
誤差は中央値 0.508\%、平均 0.832\%、最大 2.864\%
であり、弱場ほど系統的に一致が改善した。さらに \(p=60,e=0.45\)
を固定して質量比 \(m_A/m_B=1,2,4,10\) を比較すると、規格化半径 RMS
誤差は 0.508\% から 0.438\%
の範囲で安定し、放射反作用による近点縮小量は全条件でPN基準の約92\%となった。全実験で
\(|\det S_n-1|\) は浮動小数点精度内に保たれた。

本稿は一般相対論を \(S_n\)
から導出したものではない。また、状態依存の離散更新それ自体を新しい数学形式として主張するものでもない。示したのは、固定作用
\(S\) から状態依存作用 \(S_n\)
への最小拡張によって、同一の作用生成則のまま、歳差、散逸、異なる軌道形状、異なる質量比を持つ重力軌道族を内部状態更新として表現できるという実装可能性と移植性である。

\textbf{キーワード:}
状態依存写像、離散相互作用、重力軌道、Schwarzschild軌道、ポストニュートン近似、放射反作用、近点移動、思考実験

\begin{center}\rule{0.5\linewidth}{0.5pt}\end{center}

\section{0.
この稿の位置づけ}\label{ux3053ux306eux7a3fux306eux4f4dux7f6eux3065ux3051}

本稿は、研究プロジェクト設計 {[}0{]} のもとでの第四思考実験である。

第二思考実験 {[}1{]} では、内部状態

\[
X_n=\begin{pmatrix}a_n\\b_n\end{pmatrix}
\]

に対して固定した可逆線形作用 \(S\) を用い、

\[
X_{n+1}=SX_n
\]

という状態更新から、位置

\[
Y_n=\Phi(X_n)
=\begin{pmatrix}
a_n^2-b_n^2\\
2a_nb_n
\end{pmatrix}
\]

と面積時計を読み出すことで、逆二乗型の運動を生成した。構造は

\[
\boxed{
\text{内部状態}
\rightarrow
\text{相互作用}
\rightarrow
\text{軌道}
}
\]

であった。

第三思考実験 {[}2{]}
では、この向きを逆転させた。既知のニュートン二体問題と、1PN・2PN
保存項および 2.5PN
放射反作用を含むポストニュートン運動方程式から軌道を計算し、その軌道と重力波の観測データだけから、質量比・相対距離・相対時間などを読み出せるかを調べた。構造は

\[
\boxed{
\text{既知の重力モデル}
\rightarrow
\text{軌道}
\rightarrow
\text{観測}
\rightarrow
\text{状態の読出し}
}
\]

であった。

第三思考実験によって、強い場の二体軌道から複数の物理量を読み出せることは確認できた。しかし、この方法には本研究シリーズの目的から見て一つの不満が残る。

\[
\boxed{
\text{軌道そのものは、内部相互作用から生成されていない。}
}
\]

最も複雑な運動の生成を、既知のポストニュートン方程式に外部委託したままなのである。

したがって第四思考実験の問いは単純である。

\begin{center}
\fbox{\parbox{0.90\linewidth}{\centering 第三思考実験で外部計算していた重力軌道を、

\smallskip
$X_{n+1}=S(X_n)$ 型の内部相互作用へ戻せるか。}}
\end{center}

その検討から、本稿の \(S_n\) が生じた。

本稿では、札を次の意味で使う。

{\def\LTcaptype{none} % do not increment counter
\begin{longtable}[]{@{}ll@{}}
\toprule\noalign{}
札 & 意味 \\
\midrule\noalign{}
\endhead
\bottomrule\noalign{}
\endlastfoot
【原則】 & 研究系列で維持する哲学的原則 \\
【設計仮定】 & 現段階の実装のために置く仮定。将来外してよい \\
【構成】 & 既知の式を用いて \(S_n\) を明示的に作る手順 \\
【数値】 & 数値実験で確認した結果 \\
【照合】 & 既知物理との比較。基礎原理からの導出ではない \\
\end{longtable}
}

\begin{center}\rule{0.5\linewidth}{0.5pt}\end{center}

\section{1. 問い}\label{ux554fux3044}

本稿で検証する問いは次である。

\begin{center}
\fbox{\parbox{0.90\linewidth}{\centering 状態依存の逐次作用 $S_n$ を用いれば、

\smallskip
同一の作用生成則から異なる相対論的重力軌道族を生成できるか。}}
\end{center}

ここで重要なのは、個々の軌道を高精度にフィットすることではない。

一つの軌道に対して \(S_n\)
を逐次逆算すれば、その軌道を再現する写像列を構成すること自体はできる。それだけでは相互作用則を得たとは言えない。

したがって、本稿の判定基準は

\[
\boxed{
\text{同一の }\mathcal S\text{ を固定し、初期条件だけを変えて軌道族へ移植できるか}
}
\]

である。

以下ではこれを \textbf{transferability（移植性）} と呼ぶ。

\begin{center}\rule{0.5\linewidth}{0.5pt}\end{center}

\section{\texorpdfstring{2. なぜ固定作用 \(S\) から \(S_n\)
へ進むのか}{2. なぜ固定作用 S から S\_n へ進むのか}}\label{ux306aux305cux56faux5b9aux4f5cux7528-s-ux304bux3089-s_n-ux3078ux9032ux3080ux306eux304b}

\subsection{2.1
第二思考実験の固定作用}\label{ux7b2cux4e8cux601dux8003ux5b9fux9a13ux306eux56faux5b9aux4f5cux7528}

第二思考実験 {[}1{]} では、

\[
X_{n+1}=SX_n
\]

とし、同じ \(S\) を処理ごとに繰り返し作用させた。\(S\)
は可逆な線形写像であり、面積読み出しの保存から

\[
\det S=1
\]

を用いた。

この固定作用から、二次位置読み出しと面積時計を介して逆二乗型の軌道が得られた。

\subsection{\texorpdfstring{2.2 第三思考実験の軌道へ固定 \(S\)
を適用したときの問題}{2.2 第三思考実験の軌道へ固定 S を適用したときの問題}}\label{ux7b2cux4e09ux601dux8003ux5b9fux9a13ux306eux8eccux9053ux3078ux56faux5b9a-s-ux3092ux9069ux7528ux3057ux305fux3068ux304dux306eux554fux984c}

第三思考実験 {[}2{]} で扱った強い場の軌道には、単純な Kepler
型閉軌道にはない二つの特徴がある。

第一は、近点移動である。1回の動径周期に対する方位角の増分は \(2\pi\)
ではなく、半径と離心率に依存する。

第二は、2.5PN 放射反作用による軌道縮退である。処理が進むにつれて、半直弦
\(p\) と離心率 \(e\) 自体が変化する。

固定 \(S\)
のまま今回の状態表現と読出しを維持しようとすると、これらの変化を \(S\)
の外側から処理ごとに与える必要が生じる。それでは、第三思考実験で外部から軌道を与えたことと本質的に同じ問題が残る。

ここで主張しているのは、あらゆる状態空間・あらゆる読み出しを許しても固定
\(S\) では相対論的軌道を絶対に表せない、という不可能性定理ではない。

本稿で確認したのは、

\begin{center}
\fbox{\parbox{0.90\linewidth}{\centering 第二思考実験の状態表現と読み出しを保ちながら、

\smallskip
第三思考実験の軌道変化を作用の外から与えないためには、

\smallskip
$S_n=S$ という固定条件を外すのが最小である。}}
\end{center}

ということである。

\subsection{2.3
状態依存作用}\label{ux72b6ux614bux4f9dux5b58ux4f5cux7528}

そこで作用自身を、処理 \(n\) における関係状態 \(\mathcal R_n\)
の関数とする。

\[
\boxed{
S_n=\mathcal S(\mathcal R_n)
}
\]

そして

\[
\boxed{
X_{n+1}=S_nX_n
}
\]

とする。

ここで \(n\) は\textbf{物理的な時刻ではない}。

\[
\boxed{n=\text{processing step}}
\]

であり、相互作用処理の順序を表す離散番号である。物理時間 \(t\)
は後述する時計読み出しから得る。

状態依存の差分方程式や離散写像そのものは、離散力学において一般的な数学形式である
{[}3{]}。したがって、本稿は \(S\rightarrow S_n\)
という数学的一般化そのものを新規性として主張しない。

本稿での問題は、固定作用では不足した具体的な探索過程の後に、既知の一般的な離散更新形式を最小拡張として採用し、それによって同一の作用生成則から実際に重力軌道族を生成できるか、という点にある。

\begin{center}\rule{0.5\linewidth}{0.5pt}\end{center}

\section{3.
内部状態、位置読み出し、時計}\label{ux5185ux90e8ux72b6ux614bux4f4dux7f6eux8aadux307fux51faux3057ux6642ux8a08}

\subsection{3.1
二成分内部状態}\label{ux4e8cux6210ux5206ux5185ux90e8ux72b6ux614b}

内部状態を

\[
X_n=
\begin{pmatrix}
a_n\\
b_n
\end{pmatrix}
\]

とする。

観測される位置は状態そのものではなく、第二思考実験 {[}1{]}
と同じ二次読み出し

\[
\Phi(X_n)
=
\begin{pmatrix}
a_n^2-b_n^2\\
2a_nb_n
\end{pmatrix}
\]

で定義する。

したがって

\[
x_n=a_n^2-b_n^2,
\qquad
y_n=2a_nb_n,
\]

\[
r_n=a_n^2+b_n^2
\]

とすれば

\[
x_n^2+y_n^2=r_n^2.
\]

逆に

\[
X_n=
\sqrt{r_n}
\begin{pmatrix}
\cos(\phi_n/2)\\
\sin(\phi_n/2)
\end{pmatrix}
\]

と置けば

\[
\Phi(X_n)
=
\begin{pmatrix}
r_n\cos\phi_n\\
r_n\sin\phi_n
\end{pmatrix}.
\]

\subsection{3.2
処理番号と物理時間の分離}\label{ux51e6ux7406ux756aux53f7ux3068ux7269ux7406ux6642ux9593ux306eux5206ux96e2}

本稿では処理番号 \(n\) と物理時間 \(t\) を同一視しない。

\[
\boxed{n\neq t}
\]

保存系の軌道を離散化するために処理位相 \(\chi_n=n\Delta\chi\)
を用いるが、これは物理時間ではない。

第二思考実験の面積時計に対応して、数値実験02以降では

\[
\frac{dt}{d\chi}=\frac{r^2}{\sqrt p}
\]

を時計読み出しとして用いた。

したがって1処理の物理時間間隔は

\[
\Delta t_n
\simeq
\frac{r_n^2}{\sqrt{p_n}}\Delta\chi
\]

として得られる。

【設計仮定】現在の実装では \(\chi_n\) の進行幅 \(\Delta\chi\)
自体は処理順序を数値化するため外部から固定している。\(\chi\)
の増分そのものを内部関係から生成する問題は、本稿では扱わない。

\begin{center}\rule{0.5\linewidth}{0.5pt}\end{center}

\section{\texorpdfstring{4. \(\det S_n=1\)
を保つ作用の構成}{4. \textbackslash det S\_n=1 を保つ作用の構成}}\label{det-s_n1-ux3092ux4fddux3064ux4f5cux7528ux306eux69cbux6210}

軌道の半径と方位角を \(r_n,\phi_n\) とし、

\[
\theta_n=\frac{\phi_n}{2}
\]

と置く。

回転行列を

\[
R(\theta)=
\begin{pmatrix}
\cos\theta&-\sin\theta\\
\sin\theta&\cos\theta
\end{pmatrix}
\]

とする。

内部状態を含むフレームとして

\[
B_n
=
R(\theta_n)
\begin{pmatrix}
\sqrt{r_n}&0\\
0&1/\sqrt{r_n}
\end{pmatrix}
\]

を取る。

このとき

\[
S_n=B_{n+1}B_n^{-1}
\]

とすれば、

\[
S_n
=
R(\theta_{n+1})
\begin{pmatrix}
\rho_n&0\\
0&\rho_n^{-1}
\end{pmatrix}
R(-\theta_n),
\]

\[
\rho_n=
\sqrt{\frac{r_{n+1}}{r_n}}.
\]

よって

\[
\boxed{\det S_n=1}
\]

が各処理で成り立つ。

ここで注意すべきことがある。一つの状態ベクトル \(X_n\) と次状態
\(X_{n+1}\) だけから、\(2\times2\) 行列 \(S_n\)
が一意に定まるわけではない。上の \(B_n\)
は、第二のフレーム方向を逆数振幅で補って \(\det B_n=1\)
とする一つのゲージ選択である。

したがって本稿は、この \(S_n\)
が唯一の基本相互作用であるとは主張しない。

検証するのは、この単純な面積保存構成を\textbf{全ケースで変更せずに使えるか}である。

\begin{center}\rule{0.5\linewidth}{0.5pt}\end{center}

\section{\texorpdfstring{5. 保存的相対論軌道を \(S_n\)
へ載せる}{5. 保存的相対論軌道を S\_n へ載せる}}\label{ux4fddux5b58ux7684ux76f8ux5bfeux8ad6ux8eccux9053ux3092-s_n-ux3078ux8f09ux305bux308b}

\subsection{5.1 Schwarzschild
束縛軌道}\label{schwarzschild-ux675fux7e1bux8eccux9053}

Schwarzschild 時空の束縛測地線は、Darwin のパラメータ \(\chi\) を用いて
{[}4{]}

\[
r(\chi)=\frac{p}{1+e\cos\chi}
\]

と表せる。

方位角は

\[
\frac{d\phi}{d\chi}
=
\sqrt{\frac{p}{p-6-2e\cos\chi}}
\]

すなわち

\[
\frac{d\phi}{d\chi}
=
\left[
1-\frac{2(3+e\cos\chi)}{p}
\right]^{-1/2}
\]

で進む。

ここで

\[
u(\chi)=\frac{3+e\cos\chi}{p}
\]

と置くと

\[
(1-2u)^{-1/2}
=
\sum_{m=0}^{\infty}c_mu^m
\]

であり、係数は

\[
c_m=\frac{\binom{2m}{m}}{2^m}
\]

または

\[
\boxed{
c_0=1,
\qquad
c_{m+1}=\frac{2m+1}{m+1}c_m
}
\]

という漸化式で任意次数まで生成できる。

\(N\) 次で

\[
g_N(\chi;p,e)
=
\sum_{m=0}^{N}
c_m
\left(
\frac{3+e\cos\chi}{p}
\right)^m
\]

とし、

\[
\phi_{n+1}
=
\phi_n+
\int_{\chi_n}^{\chi_{n+1}}
g_N(\chi;p,e)d\chi
\]

と更新する。

本数値実験では

\[
N=12
\]

を用いた。

\subsection{5.2
この構成の意味}\label{ux3053ux306eux69cbux6210ux306eux610fux5473}

この段階では、\(r(\chi)\) と \(\phi(\chi)\) を既知の Schwarzschild
軌道から構成し、それを \(S_n\) に持ち上げている。

したがって、\(S_n\) がこの構成元の Schwarzschild
軌道そのものを再現することは、独立な物理検証ではない。

これは

\[
\boxed{
\text{相対論的軌道を }\det S_n=1\text{ の逐次作用として表現できる}
}
\]

ことを示す\textbf{構成}である。

非自明な試験は、その後、この同一構成を変更せずに、第三思考実験で独立に生成した有限質量比PN軌道と比較し、初期条件を変えても移植できるかを見ることである。

\begin{center}\rule{0.5\linewidth}{0.5pt}\end{center}

\section{\texorpdfstring{6. 放射反作用を含む \(p_n,e_n\)
の漸化}{6. 放射反作用を含む p\_n,e\_n の漸化}}\label{ux653eux5c04ux53cdux4f5cux7528ux3092ux542bux3080-p_ne_n-ux306eux6f38ux5316}

\subsection{6.1
基準となる2.5PN放射反作用}\label{ux57faux6e96ux3068ux306aux308b2.5pnux653eux5c04ux53cdux4f5cux7528}

第三思考実験 {[}2{]} の基準軌道は、調和座標で Newtonian + 1PN + 2PN
保存項、および 2.5PN 放射反作用を用いている。2PN
までは保存的であり、最初の非保存的効果である放射反作用が 2.5PN
に現れることは標準的なPN理論である {[}5,7{]}。

本稿では第三思考実験と同じ Kidder {[}5{]} の非自転部分を使用した。

\[
\mathbf a_{\rm RR}
=C(A\mathbf n+B\mathbf v),
\]

\[
C=\frac85\nu r^{-3},
\]

\[
A
=\dot r
\left(
18v^2+\frac{2}{3r}-25\dot r^2
\right),
\]

\[
B
=-\left(
6v^2-\frac2r-15\dot r^2
\right),
\]

ここで

\[
\nu=\frac{m_Am_B}{(m_A+m_B)^2}
\]

である。

\subsection{6.2 最小の osculating-element
更新}\label{ux6700ux5c0fux306e-osculating-element-ux66f4ux65b0}

放射反作用を \(S_n\) 系へ加えるため、現在の軌道を瞬間的な Newtonian
osculating orbit として読み、

\[
E=\frac{v^2}{2}-\frac1r,
\]

\[
h=(\mathbf r\times\mathbf v)_z,
\]

\[
p=h^2,
\]

\[
e^2=1+2Ep
\]

を用いた。

放射反作用による変化は

\[
\dot E
=\mathbf v\cdot\mathbf a_{\rm RR}
=C(A\dot r+Bv^2),
\]

また \(\mathbf r\times\mathbf n=0\) より

\[
\dot h=CBh,
\]

したがって

\[
\boxed{
\dot p=2CBp
}
\]

となる。

さらに

\[
e^2=1+2Ep
\]

を微分して

\[
\boxed{
\dot e
=\frac{p\dot E+E\dot p}{e}
}
\]

を得る。

実装では \(dt/d\chi=r^2/\sqrt p\) を用いて、\(p(\chi),e(\chi)\) を4次
Runge--Kutta 法で更新した。

軌道要素を局所的に変化させて加速された軌道を osculating orbit
の列として扱う考え方自体は既知であり、Schwarzschild
時空に対しても体系化されている
{[}6{]}。したがって、この軌道要素更新そのものを本稿の新規性とはしない。

本稿の試験は、それを内部状態の

\[
X_n\rightarrow S_n\rightarrow X_{n+1}
\]

という面積保存作用列へ組み込み、同じ \(S_n\)
生成則を複数軌道へ移植できるかにある。

\begin{center}\rule{0.5\linewidth}{0.5pt}\end{center}

\section{7.
数値実験の設計}\label{ux6570ux5024ux5b9fux9a13ux306eux8a2dux8a08}

\subsection{7.1 共通条件}\label{ux5171ux901aux6761ux4ef6}

単位は第三思考実験と同じく

\[
G=c=M=1,
\qquad
M=m_A+m_B
\]

とした。

\(S_n\) の保存部分の展開次数は全ケースで

\[
N=12
\]

とした。

ケースごとに \(S_n\) の係数を再フィットする操作は行っていない。

また全ケースで

\[
\det S_n=1
\]

を数値的に検査した。

質量比について、本稿では第二思考実験の面積保存量との記号衝突を避けるため

\[
q_m=\frac{m_A}{m_B}\ge1
\]

と書く。保存されたCSVの列名 \texttt{q} はこの \(q_m\) を意味する。

\subsection{\texorpdfstring{7.2 数値実験01：保存的
\(S_n\)}{7.2 数値実験01：保存的 S\_n}}\label{ux6570ux5024ux5b9fux9a1301ux4fddux5b58ux7684-s_n}

第三思考実験の C1 と同じ初期条件

\[
a=30,
\qquad b=27,
\qquad q_m=1
\]

を使った。

\[
e_0
=\sqrt{1-\frac{b^2}{a^2}}
=0.435889894\ldots,
\]

\[
p_0=a(1-e_0^2)=\frac{b^2}{a}=24.3.
\]

\(p,e\) は固定し、保存的 \(S_n\) だけで3動径周期を生成した。

\subsection{\texorpdfstring{7.3 数値実験02：\(S_n\) +
2.5PN放射反作用}{7.3 数値実験02：S\_n + 2.5PN放射反作用}}\label{ux6570ux5024ux5b9fux9a1302s_n-2.5pnux653eux5c04ux53cdux4f5cux7528}

数値実験01と同じ初期条件に、§6 の \(p_n,e_n\) 更新を追加した。

\(S_n\) の行列構造は変更していない。

\subsection{7.4 数値実験03：9軌道
transferability}\label{ux6570ux5024ux5b9fux9a13039ux8eccux9053-transferability}

次に質量比を

\[
q_m=1
\]

に固定し、

\[
p\in\{24.3,60,120\},
\]

\[
e\in\{0.20,0.45,0.55\}
\]

の直積による9条件へ、数値実験02と\textbf{全く同じ \(S_n\)
生成アルゴリズム}を適用した。

当初は高離心率条件として \(e=0.65\) を予定したが、最強場
\(p=24.3,e=0.65\)
では比較基準である第三思考実験側のPN初期速度較正器が安定に解を生成できなかった。このため、\(S_n\)
側だけに有利な比較領域を選ばないよう、共通に安定して生成できる
\(e=0.55\) を高離心率条件とした。

\subsection{7.5 数値実験04：質量比
transferability}\label{ux6570ux5024ux5b9fux9a1304ux8ceaux91cfux6bd4-transferability}

さらに

\[
p=60,
\qquad e=0.45
\]

を固定し、

\[
q_m\in\{1,2,4,10\}
\]

とした。

\(q_m=1\) は数値実験03と重なるため、追加計算として新しい質量比は
\(2,4,10\) の3条件である。

\begin{center}\rule{0.5\linewidth}{0.5pt}\end{center}

\section{8. 数値結果}\label{ux6570ux5024ux7d50ux679c}

\subsection{8.1
数値実験01・02：C1基準軌道}\label{ux6570ux5024ux5b9fux9a130102c1ux57faux6e96ux8eccux9053}

第三思考実験のPN基準では、抽出された近点移動は

\[
\Delta\varpi_{\rm PN}=0.926885\ {\rm rad/orbit}
\]

であり、近点半径は

\[
16.923303
\rightarrow
16.825431
\rightarrow
16.725070
\]

と縮小した。

保存的 \(S_n\) では

\[
\Delta\varpi_{S,\rm cons}=0.960214\ {\rm rad/orbit},
\]

\[
\frac{\mathrm{RMS}_r}{a}=0.0251045
\]

であり、放射反作用を持たないため近点半径は \(16.923303\)
のままであった。

2.5PN 漸化を加えると、3周で

\[
p:
24.3\rightarrow23.704787,
\]

\[
e:
0.4358899\rightarrow0.4203470
\]

と変化し、

\[
\frac{\mathrm{RMS}_r}{a}
=0.0195011
\]

となった。

これは保存系の 2.510\% から 1.950\% への低下、すなわち約 22.3\%
の改善である。

一方、近点移動は

\[
\Delta\varpi_{S,\rm RR}=0.969957\ {\rm rad/orbit}
\]

となり、PN基準よりやや大きい。

これは、散逸を入れることで半径方向は改善する一方、保存部分が
Schwarzschild/test-particle
型で、基準PNが有限質量比の1PN+2PN保存補正を含む差が残っていることと整合的である。【照合】

\begin{figure}
\centering
\pandocbounded{\includegraphics[keepaspectratio,alt={図1　第三思考実験PN基準、保存的S\_n、2.5PN放射反作用を加えたS\_nの比較}]{figs/compare_C1_PN_vs_Exp01_Exp02_N12.png}}
\caption{図1　第三思考実験PN基準、保存的S\_n、2.5PN放射反作用を加えたS\_nの比較}
\end{figure}

\textbf{図1.} C1 初期条件における3系の比較。保存的 \(S_n\)
に2.5PN放射反作用漸化を追加すると、PN基準に見られる軌道縮退が現れ、半径方向の一致が改善する。

\subsection{8.2
数値実験03：9軌道}\label{ux6570ux5024ux5b9fux9a13039ux8eccux9053}

表1に9条件の結果を示す。

\subsubsection{\texorpdfstring{表1　同一 \(S_n\)
生成則の9軌道への移植}{表1　同一 S\_n 生成則の9軌道への移植}}\label{ux88681-ux540cux4e00-s_n-ux751fux6210ux5247ux306e9ux8eccux9053ux3078ux306eux79fbux690d}

{\def\LTcaptype{none} % do not increment counter
\begin{longtable}[]{@{}
  >{\raggedleft\arraybackslash}p{(\linewidth - 10\tabcolsep) * \real{0.1667}}
  >{\raggedleft\arraybackslash}p{(\linewidth - 10\tabcolsep) * \real{0.1667}}
  >{\raggedleft\arraybackslash}p{(\linewidth - 10\tabcolsep) * \real{0.1667}}
  >{\raggedleft\arraybackslash}p{(\linewidth - 10\tabcolsep) * \real{0.1667}}
  >{\raggedleft\arraybackslash}p{(\linewidth - 10\tabcolsep) * \real{0.1667}}
  >{\raggedleft\arraybackslash}p{(\linewidth - 10\tabcolsep) * \real{0.1667}}@{}}
\toprule\noalign{}
\begin{minipage}[b]{\linewidth}\raggedleft
\(p\)
\end{minipage} & \begin{minipage}[b]{\linewidth}\raggedleft
\(e\)
\end{minipage} & \begin{minipage}[b]{\linewidth}\raggedleft
radial RMS/\(a\)
\end{minipage} & \begin{minipage}[b]{\linewidth}\raggedleft
PN近点移動 {[}rad/orbit{]}
\end{minipage} & \begin{minipage}[b]{\linewidth}\raggedleft
\(S_n\)近点移動 {[}rad/orbit{]}
\end{minipage} & \begin{minipage}[b]{\linewidth}\raggedleft
差 {[}rad/orbit{]}
\end{minipage} \\
\midrule\noalign{}
\endhead
\bottomrule\noalign{}
\endlastfoot
24.3 & 0.20 & 0.783\% & 0.927796 & 0.966421 & +0.038625 \\
24.3 & 0.45 & 2.045\% & 0.926007 & 0.970253 & +0.044246 \\
24.3 & 0.55 & 2.864\% & 0.926956 & 0.972619 & +0.045663 \\
60 & 0.20 & 0.188\% & 0.332836 & 0.340217 & +0.007381 \\
60 & 0.45 & 0.508\% & 0.331147 & 0.340531 & +0.009385 \\
60 & 0.55 & 0.692\% & 0.331805 & 0.340725 & +0.008919 \\
120 & 0.20 & 0.051\% & 0.160726 & 0.163264 & +0.002538 \\
120 & 0.45 & 0.150\% & 0.161215 & 0.163328 & +0.002112 \\
120 & 0.55 & 0.208\% & 0.161834 & 0.163367 & +0.001533 \\
\end{longtable}
}

9ケース全体では

\[
\operatorname{median}
\left(\frac{\mathrm{RMS}_r}{a}\right)
=0.508392\%,
\]

\[
\operatorname{mean}
\left(\frac{\mathrm{RMS}_r}{a}\right)
=0.832314\%,
\]

\[
\max
\left(\frac{\mathrm{RMS}_r}{a}\right)
=2.8642\%.
\]

近点移動誤差の絶対値は中央値 \(0.00891949\) rad/orbit、平均
\(0.0178224\) rad/orbit、最大 \(0.0456626\) rad/orbit であった。

\begin{figure}
\centering
\pandocbounded{\includegraphics[keepaspectratio,alt={図2　9軌道条件におけるPN基準とS\_n軌道の比較}]{figs/experiment03_9case_orbit_grid.png}}
\caption{図2　9軌道条件におけるPN基準とS\_n軌道の比較}
\end{figure}

\textbf{図2.} \(p\) と \(e\)
を変えた9条件。ケース固有の係数フィットを行わず、同じ \(S_n\)
生成則を使っている。

\begin{figure}
\centering
\pandocbounded{\includegraphics[keepaspectratio,alt={図3　9軌道条件の規格化半径RMS誤差}]{figs/experiment03_radial_rms_over_a.png}}
\caption{図3　9軌道条件の規格化半径RMS誤差}
\end{figure}

\textbf{図3.} 規格化半径 RMS
誤差。弱場側ほど誤差が系統的に減少し、\(p=120\) では全3条件で 0.21\%
以下となった。

近点縮小についても全9条件でPN基準と同じ符号・同じ傾向を再現した。ただし
\(S_n\)
側は縮小量を系統的に小さく出し、その相対差は強場で約22\%、中間で約8\%、弱場で約4\%であった。

重要なのは、誤差がケースごとに無秩序に発散せず、場の強さと離心率に対して滑らかに変化したことである。

\subsection{8.3
数値実験04：質量比}\label{ux6570ux5024ux5b9fux9a1304ux8ceaux91cfux6bd4}

表2に \(p=60,e=0.45\) で質量比のみを変えた結果を示す。

\subsubsection{表2　質量比を変えた
transferability}\label{ux88682-ux8ceaux91cfux6bd4ux3092ux5909ux3048ux305f-transferability}

{\def\LTcaptype{none} % do not increment counter
\begin{longtable}[]{@{}
  >{\raggedleft\arraybackslash}p{(\linewidth - 10\tabcolsep) * \real{0.1667}}
  >{\raggedleft\arraybackslash}p{(\linewidth - 10\tabcolsep) * \real{0.1667}}
  >{\raggedleft\arraybackslash}p{(\linewidth - 10\tabcolsep) * \real{0.1667}}
  >{\raggedleft\arraybackslash}p{(\linewidth - 10\tabcolsep) * \real{0.1667}}
  >{\raggedleft\arraybackslash}p{(\linewidth - 10\tabcolsep) * \real{0.1667}}
  >{\raggedleft\arraybackslash}p{(\linewidth - 10\tabcolsep) * \real{0.1667}}@{}}
\toprule\noalign{}
\begin{minipage}[b]{\linewidth}\raggedleft
\(q_m=m_A/m_B\)
\end{minipage} & \begin{minipage}[b]{\linewidth}\raggedleft
\(\nu\)
\end{minipage} & \begin{minipage}[b]{\linewidth}\raggedleft
radial RMS/\(a\)
\end{minipage} & \begin{minipage}[b]{\linewidth}\raggedleft
PN近点移動 {[}rad/orbit{]}
\end{minipage} & \begin{minipage}[b]{\linewidth}\raggedleft
\(S_n\)近点移動 {[}rad/orbit{]}
\end{minipage} & \begin{minipage}[b]{\linewidth}\raggedleft
差 {[}rad/orbit{]}
\end{minipage} \\
\midrule\noalign{}
\endhead
\bottomrule\noalign{}
\endlastfoot
1 & 0.25000000 & 0.5084\% & 0.331147 & 0.340531 & 0.009385 \\
2 & 0.22222222 & 0.4967\% & 0.332513 & 0.340496 & 0.007983 \\
4 & 0.16000000 & 0.4704\% & 0.331694 & 0.340419 & 0.008725 \\
10 & 0.08264463 & 0.4379\% & 0.333864 & 0.340322 & 0.006458 \\
\end{longtable}
}

規格化半径誤差は

\[
0.5084\%
\rightarrow
0.4967\%
\rightarrow
0.4704\%
\rightarrow
0.4379\%
\]

と、質量比を大きくしても増大しなかった。

\begin{figure}
\centering
\pandocbounded{\includegraphics[keepaspectratio,alt={図4　質量比4条件のPN基準軌道とS\_n軌道}]{figs/experiment04_mass_ratio_orbits.png}}
\caption{図4　質量比4条件のPN基準軌道とS\_n軌道}
\end{figure}

\textbf{図4.} \(p=60,e=0.45\) を固定し、\(q_m=1,2,4,10\)
とした比較。同一の \(S_n\) 生成則を使用した。

\begin{figure}
\centering
\pandocbounded{\includegraphics[keepaspectratio,alt={図5　質量比に対する規格化半径RMS誤差}]{figs/experiment04_rms_vs_q.png}}
\caption{図5　質量比に対する規格化半径RMS誤差}
\end{figure}

\textbf{図5.}
質量比を変えても誤差は安定し、極端質量比側でわずかに低下する。

放射反作用による近点縮小量の比

\[
\frac{\Delta r_{p,S}}{\Delta r_{p,\rm PN}}
\]

は、\(q_m=1,2,4,10\) に対してそれぞれ約

\[
0.9178,
\quad
0.9177,
\quad
0.9185,
\quad
0.9191
\]

であった。

\begin{figure}
\centering
\pandocbounded{\includegraphics[keepaspectratio,alt={図6　対称質量比nuに対する近点縮小量}]{figs/experiment04_shrink_vs_nu.png}}
\caption{図6　対称質量比nuに対する近点縮小量}
\end{figure}

\textbf{図6.} \(S_n\)
の近点縮小量は、質量比を変えてもPN基準の約92\%というほぼ一定の比率を保つ。

これは、現在の最小漸化が質量比依存を無秩序に失っているのではなく、約8\%の系統的な補正が不足している可能性を示す。【照合】

\subsection{8.4 面積保存}\label{ux9762ux7a4dux4fddux5b58}

全9軌道試験で

\[
\max |\det S_n-1|
\le6.661\times10^{-16},
\]

質量比試験でも

\[
\max |\det S_n-1|
\lesssim5.6\times10^{-16}
\]

であった。

したがって、物理空間で読み出される軌道が放射反作用によって縮退している間も、内部2状態空間の各
\(S_n\) は数値精度内で面積保存を維持した。【数値】

\begin{center}\rule{0.5\linewidth}{0.5pt}\end{center}

\section{9.
何が確認できたか}\label{ux4f55ux304cux78baux8a8dux3067ux304dux305fux304b}

\subsection{9.1
一つの軌道を作れることではなく、軌道族へ移せること}\label{ux4e00ux3064ux306eux8eccux9053ux3092ux4f5cux308cux308bux3053ux3068ux3067ux306fux306aux304fux8eccux9053ux65cfux3078ux79fbux305bux308bux3053ux3068}

本稿で最も重要なのは、C1 一例の一致精度ではない。

一つの既知軌道から \(S_n\)
を逆算するだけなら、それは軌道の符号化にすぎない。

そこで実験03では、\(S_n\) の生成則を固定したまま \(p,e\)
の9条件へ適用した。さらに実験04では質量比も変えた。

その結果、

\[
\boxed{
\text{同じ }\mathcal S\text{ が、場の強さ・離心率・質量比の異なる軌道族で機能した。}
}
\]

この意味で、今回の結果は単一軌道への case-specific fitting より強い。

\subsection{9.2
弱場・極端質量比で誤差が減る}\label{ux5f31ux5834ux6975ux7aefux8ceaux91cfux6bd4ux3067ux8aa4ux5deeux304cux6e1bux308b}

保存部分は Schwarzschild/test-particle 軌道から構成している。

そのため、

\[
p\rightarrow\infty
\]

の弱場側、および

\[
\nu\rightarrow0
\]

の極端質量比側で、有限質量比PN基準との差が減ることは自然な傾向である。

実際に、9軌道試験では \(p\)
の増加とともに誤差が系統的に低下し、質量比試験でも \(q_m\)
の増加、すなわち \(\nu\) の減少とともに誤差がわずかに低下した。

これは結果の解釈であり、基礎原理からの導出ではない。

\subsection{9.3
内部保存と観測上の散逸は同じではない}\label{ux5185ux90e8ux4fddux5b58ux3068ux89b3ux6e2cux4e0aux306eux6563ux9038ux306fux540cux3058ux3067ux306fux306aux3044}

各処理で

\[
\det S_n=1
\]

である一方、読み出された軌道では \(p_n,e_n\)
が減少し、近点・遠点が縮小する。

したがって今回の構成では、

\[
\boxed{
\text{内部状態更新の面積保存}
\neq
\text{観測された軌道の非散逸性}
}
\]

である。

散逸を \(\det S_n<1\)
のような内部面積の消失として直接実装しなくても、状態依存する作用列と読み出しの組によって、散逸的な有効運動を表現できる。

ただし、これを自然界の放射反作用の基礎機構と同一視することは本稿の範囲を超える。

\begin{center}\rule{0.5\linewidth}{0.5pt}\end{center}

\section{10.
既知の数学・物理との関係}\label{ux65e2ux77e5ux306eux6570ux5b66ux7269ux7406ux3068ux306eux95a2ux4fc2}

\subsection{10.1
状態依存の離散更新}\label{ux72b6ux614bux4f9dux5b58ux306eux96e2ux6563ux66f4ux65b0}

差分方程式・離散力学では、次状態を現在状態の関数として

\[
X_{n+1}=F_n(X_n)
\]

あるいは自律系として

\[
X_{n+1}=F(X_n)
\]

と書くことは一般的である {[}3{]}。

したがって、作用を

\[
S_n=\mathcal S(\mathcal R_n)
\]

と状態依存にしたこと自体は、本稿だけのために導入した特殊な数学操作ではない。

本稿固有の探索上の意味は、第二思考実験の固定作用を実際に第三思考実験の相対論的軌道へ拡張しようとして不足を認識し、\textbf{状態や読み出しを増やす前に、固定作用という一つの仮定だけを緩めた}ことにある。

\subsection{10.2 Schwarzschild 軌道}\label{schwarzschild-ux8eccux9053}

保存部分で用いた \(p,e,\chi\) による束縛軌道の表現は Schwarzschild
測地線の既知の構造であり、Darwin に遡る {[}4{]}。

本稿はこの測地線解を新たに導出したものではなく、その軌道進行を
\(\det S_n=1\) の状態更新列へ持ち上げている。

\subsection{10.3
PN基準と放射反作用}\label{pnux57faux6e96ux3068ux653eux5c04ux53cdux4f5cux7528}

比較に用いた有限質量比軌道と2.5PN放射反作用は、第三思考実験 {[}2{]}
と同じ既知のポストニュートン力学であり
{[}5,7{]}、一般相対論の代替として本稿が発明した項ではない。

2PN までの保存力学と 2.5PN
で現れる最初の放射反作用という次数構造も標準的である {[}7{]}。

\subsection{10.4 osculating orbit}\label{osculating-orbit}

加速された軌道を瞬間ごとの軌道要素の変化として扱う方法も既知であり、Schwarzschild
時空での osculating-orbit 法として整理されている {[}6{]}。

したがって本稿が新規性として問うのは、これら既知の物理式の存在ではなく、

\begin{center}
\fbox{\parbox{0.90\linewidth}{\centering それらの軌道変化を、同一の内部作用形式 $X_{n+1}=S_nX_n$ へ載せられるか}}
\end{center}

という構成上の問いである。

\begin{center}\rule{0.5\linewidth}{0.5pt}\end{center}

\section{11. 限界}\label{ux9650ux754c}

本稿の結果には明確な限界がある。

\begin{enumerate}
\def\labelenumi{\arabic{enumi}.}
\item
  \textbf{一般相対論を \(S_n\) から導出していない。}\\
  保存部分は Schwarzschild
  軌道を用いて構成し、散逸部分は既知の2.5PN放射反作用を用いている。本稿は表現可能性・実装可能性の試験である。
\item
  \textbf{\(S_n\) は一意ではない。}\\
  \(X_n\rightarrow X_{n+1}\) だけでは \(2\times2\)
  行列は一意に定まらない。本文の \(B_n\) は \(\det S_n=1\)
  を維持する一つの自然な補完である。
\item
  \textbf{実験02以降は hybrid 構成である。}\\
  保存角度は Schwarzschild/test-particle
  型であり、散逸項は有限質量比の調和座標2.5PN式である。両者の座標・近似次数を完全に統一した理論ではない。
\item
  \textbf{\(p,e\) 更新は最小の Newtonian osculating relations
  を使う。}\\
  PN補正済みの軌道要素変換まで精密化していない。このため約8\%の系統的な縮退差や強場での歳差差が残る可能性がある。
\item
  \textbf{処理位相 \(\chi\) の増分は内生化していない。}\\
  \(n\) は時間ではなく処理番号であり、\(t\)
  は読み出すが、現在のプログラムでは \(\Delta\chi\)
  を数値処理の刻みとして与えている。処理順序そのものを関係状態だけから生成する問題は残る。
\item
  \textbf{比較対象自身も完全な数値相対論ではない。}\\
  第三思考実験の基準は Newtonian + 1PN + 2PN + 2.5PN
  の有限次数PN近似である。したがって本稿の誤差は完全な一般相対論との差ではなく、この基準モデルとの差である。
\item
  \textbf{固定 \(S\) の一般的不可能性は証明していない。}\\
  本稿で示したのは、第二思考実験の状態表現と読み出しを保つ今回の構成では、固定
  \(S\) の外に処理依存情報を置かずに第三思考実験の軌道変化を実装するには
  \(S_n\) への一般化が最小であった、という探索結果である。
\end{enumerate}

これらの限界は、本稿の主張を

\[
\boxed{
S_n\text{ が重力の基礎法則である}
}
\]

ではなく、

\[
\boxed{
S_n\text{ 型の逐次相互作用が、重力軌道族を載せられる十分な表現能力を持つ}
}
\]

に限定する理由である。

\begin{center}\rule{0.5\linewidth}{0.5pt}\end{center}

\section{12. 再現性}\label{ux518dux73feux6027}

本稿の全数値実験は、本文と同じ Google Drive
プロジェクトフォルダに保存した。

\subsection{12.1 数値実験01・02}\label{ux6570ux5024ux5b9fux9a130102}

主なファイル：

\begin{itemize}
\tightlist
\item
  \texttt{run\_experiment01\_sn\_conservative.py}
\item
  \texttt{run\_experiment02\_sn\_2p5pn\_rr.py}
\item
  \texttt{compare\_experiments\_01\_02.py}
\item
  \texttt{run\_all\_experiments01\_02.py}
\item
  \texttt{reference\_pn\_C1.csv}
\item
  \texttt{experiment01\_sn\_conservative\_C1\_N12.csv}
\item
  \texttt{experiment02\_sn\_2p5pn\_rr\_C1\_N12.csv}
\item
  \texttt{comparison\_experiments\_01\_02\_C1\_N12.json}
\item
  \texttt{compare\_C1\_PN\_vs\_Exp01\_Exp02\_N12.svg}
\end{itemize}

\subsection{12.2 数値実験03}\label{ux6570ux5024ux5b9fux9a1303}

\texttt{experiment03\_9case\_transferability/} に保存した。

各ケースについて、図化前の raw data として

\begin{itemize}
\tightlist
\item
  \texttt{pn\_integrator\_*.csv}：PN RK4 積分の生row
\item
  \texttt{pn\_plot\_*.csv}：比較図に用いた一様時刻PN row
\item
  \texttt{sn\_*.csv}：\(S_n\) の全処理step row
\end{itemize}

を保存した。9条件なので raw CSV は計27本である。

図化コードは保存済みCSVを読むだけで、力学を再計算しない。

\subsection{12.3 数値実験04}\label{ux6570ux5024ux5b9fux9a1304}

\texttt{experiment04\_mass\_ratio\_transferability/} に保存した。

\(q_m=1,2,4,10\) の4条件について同じ3種類の raw CSV、計12本を保存した。

すべての実験フォルダに、集計CSV/JSON、実行ログ、SHA-256 checksum
を保存している。

したがって、別の誤差指標や追加解析は、軌道生成計算を再実行せずに raw row
から行える。

\begin{center}\rule{0.5\linewidth}{0.5pt}\end{center}

\section{13. 結論}\label{ux7d50ux8ad6}

第二思考実験では、固定作用

\[
X_{n+1}=SX_n
\]

から、積の位置読み出しと面積時計を通して逆二乗型の軌道を生成した。

第三思考実験では方向を逆転させ、既知の重力方程式によって生成した軌道から、質量比・相対距離・相対時間を読み出した。

しかしその段階では、軌道そのものが外部の既知理論から与えられていた。

本稿ではこの不満を出発点として、第三思考実験の相対論的軌道を再び内部相互作用から生成できないかを検討した。

固定 \(S\)
のままでは、今回の状態表現と読出しを保ちながら近点移動と軌道縮退に必要な処理依存性を作用の内部に持たせることができないため、固定作用という制約だけを外し、

\[
\boxed{
X_{n+1}=S_nX_n,
\qquad
S_n=\mathcal S(\mathcal R_n)
}
\]

とした。

\(n\) は物理時刻ではなく処理番号であり、時間は面積時計から別途読み出す。

この \(S_n\) を \(\det S_n=1\) の回転と reciprocal squeeze
から構成し、Schwarzschild 型の保存軌道と、2.5PN 放射反作用による
\(p_n,e_n\) の変化を組み込んだ。

その結果、単一の強場軌道だけでなく、

\[
p=\{24.3,60,120\},
\qquad
e=\{0.20,0.45,0.55\}
\]

の9軌道、および

\[
q_m=\{1,2,4,10\}
\]

の異なる質量比に、同一の作用生成則をケース固有の再フィットなしで適用できた。

9軌道の規格化半径 RMS 誤差は中央値 0.508\%、平均 0.832\%、最大 2.864\%
で、弱場ほど系統的に小さくなった。質量比試験では 0.508\% から 0.438\%
の範囲で安定し、近点縮小量はPN基準の約92\%というほぼ一定比率を保った。

したがって本稿で確認されたのは、

\[
\boxed{
\text{一つの既知軌道を }S_n\text{ で符号化できること}
}
\]

だけではない。

より重要なのは、

\[
\boxed{
\text{同一の状態依存作用生成則が、異なる重力軌道族へ移植できること}
}
\]

である。

これは \(S_n\)
が一般相対論から独立に導出されたことを意味しない。しかし、研究系列の元来の問いである

\[
\boxed{
\text{物理運動を、状態と関係の逐次相互作用から生成できるか}
}
\]

に対して、固定 \(S\) から状態依存 \(S_n\)
への最小拡張が有効な実装候補であることを示している。

次の問題は、今回 Schwarzschild 軌道と2.5PN式から構成した
\(\mathcal S(\mathcal R)\)
自体を、さらに少数の関係的不変量から導出できるかである。

\begin{center}\rule{0.5\linewidth}{0.5pt}\end{center}

\section{参考文献}\label{ux53c2ux8003ux6587ux732e}

\subsection{本系列}\label{ux672cux7cfbux5217}

{[}0{]} 木原範昭,
\textbf{状態と関係性から時空・粒子・相互作用の創発を探索する研究プロジェクト設計
― 不確実性下の基礎物理モデル探索に対する予備設計と実装可能性検証の枠組み
―}, v1.0 (2026-09-20). Concept DOI: 10.5281/zenodo.22851944; Version
DOI: 10.5281/zenodo.22851945.

{[}1{]} 木原範昭,
\textbf{第二思考実験：クーロン型の逆二乗の力を、積の読み出しと面積の時計から導く
―
二つの値と線形の法則のまま、焦点・引力と斥力・中性を出し、二つの値では書けないものを確定する
―}, v1.2 (2026-09-21). Concept DOI: 10.5281/zenodo.22867336; Version
DOI: 10.5281/zenodo.22876596.

{[}2{]} 木原範昭, \textbf{第三思考実験：a b 2体の天体の軌道から a b
2体の状態を読み出す ― 意味を定めない二つの値と G=c=M=1
の知識だけで、軌道と重力波の読み出しから質量比・相対距離・相対時間を得る
―}, v1.5 (2026-09-23). Concept DOI: 10.5281/zenodo.22909660; Version
DOI: 10.5281/zenodo.22909661.

\subsection{外部}\label{ux5916ux90e8}

{[}3{]} S. Elaydi, \textbf{An Introduction to Difference Equations}, 3rd
ed., Undergraduate Texts in Mathematics, Springer, New York (2005). DOI:
10.1007/0-387-27602-5.

{[}4{]} C. G. Darwin, ``The Gravity Field of a Particle,''
\emph{Proceedings of the Royal Society of London. Series A, Mathematical
and Physical Sciences}, \textbf{249}(1257), 180--194 (1959). DOI:
10.1098/rspa.1959.0015.

{[}5{]} L. E. Kidder, ``Coalescing binary systems of compact objects to
(post)\(^{5/2}\)-Newtonian order. V. Spin effects,'' \emph{Physical
Review D}, \textbf{52}, 821--847 (1995). DOI: 10.1103/PhysRevD.52.821.

{[}6{]} A. Pound and E. Poisson, ``Osculating orbits in Schwarzschild
spacetime, with an application to extreme mass-ratio inspirals,''
\emph{Physical Review D}, \textbf{77}, 044013 (2008). DOI:
10.1103/PhysRevD.77.044013.

{[}7{]} L. Blanchet, ``Post-Newtonian theory for gravitational waves,''
\emph{Living Reviews in Relativity}, \textbf{27}, Article 4 (2024). DOI:
10.1007/s41114-024-00050-z.

\begin{center}\rule{0.5\linewidth}{0.5pt}\end{center}

\section{付録A：実験ファイルの相対パス}\label{ux4ed8ux9332aux5b9fux9a13ux30d5ux30a1ux30a4ux30ebux306eux76f8ux5bfeux30d1ux30b9}

本稿を保存するフォルダを基準として、本文で使用した図は次のとおりである。

\begin{enumerate}
\def\labelenumi{\arabic{enumi}.}
\tightlist
\item
  \texttt{compare\_C1\_PN\_vs\_Exp01\_Exp02\_N12.svg}
\item
  \texttt{experiment03\_9case\_transferability/figures/experiment03\_9case\_orbit\_grid.svg}
\item
  \texttt{experiment03\_9case\_transferability/figures/experiment03\_radial\_rms\_over\_a.svg}
\item
  \texttt{experiment04\_mass\_ratio\_transferability/figures/experiment04\_mass\_ratio\_orbits.svg}
\item
  \texttt{experiment04\_mass\_ratio\_transferability/figures/experiment04\_rms\_vs\_q.svg}
\item
  \texttt{experiment04\_mass\_ratio\_transferability/figures/experiment04\_shrink\_vs\_nu.svg}
\end{enumerate}

本稿の表1・表2は、それぞれ次の保存済み集計データから転記した。

\begin{itemize}
\tightlist
\item
  \texttt{experiment03\_9case\_transferability/summary/case\_metrics.csv}
\item
  \texttt{experiment04\_mass\_ratio\_transferability/summary/case\_metrics.csv}
\end{itemize}

\end{document}
